% SPDX-License-Identifier: MPL-2.0
% Copyright (c) 2026 Jonathan D.A. Jewell <j.d.a.jewell@open.ac.uk>
%
% SNIFs: Safer Native Implemented Functions for the BEAM via WebAssembly Sandboxing
% Prepared for Zenodo deposit.

\documentclass[11pt,a4paper]{article}

%%% Packages %%%
\usepackage[utf8]{inputenc}
\usepackage[T1]{fontenc}
\usepackage{lmodern}
\usepackage{amsmath,amssymb,amsthm}
\usepackage{mathtools}
\usepackage{listings}
\usepackage{booktabs}
\usepackage{tabularx}
\usepackage{hyperref}
\usepackage[margin=2.5cm]{geometry}
\usepackage{xcolor}
\usepackage{enumitem}
\usepackage{microtype}

%%% Theorem environments %%%
\newtheorem{definition}{Definition}
\newtheorem{theorem}{Theorem}
\newtheorem{observation}{Observation}

%%% Listings %%%
\definecolor{codegreen}{rgb}{0,0.5,0}
\definecolor{codegray}{rgb}{0.5,0.5,0.5}
\definecolor{codepurple}{rgb}{0.58,0,0.82}
\definecolor{backcolour}{rgb}{0.97,0.97,0.97}

\lstdefinelanguage{Zig}{
  keywords={const,var,fn,export,pub,return,if,else,while,for,unreachable,
            comptime,inline,struct,enum,union,error,try,catch,defer,
            null,undefined,true,false,void,noreturn,anytype},
  sensitive=true,
  morecomment=[l]{//},
  morestring=[b]",
}

\lstdefinelanguage{Elixir}{
  keywords={def,defmodule,do,end,case,with,when,fn,if,else,
            true,false,nil,and,or,not,in},
  sensitive=true,
  morecomment=[l]{\#},
  morestring=[b]",
}

\lstdefinelanguage{WebAssembly}{
  keywords={unreachable,local.get,local.set,i32.const,i64.const,
            call,br\_if,block,end,return,i32.add,memory.fill,
            i32.store,i64.store,global.get,global.set},
  sensitive=true,
  morecomment=[l]{;; },
}

\lstset{
  backgroundcolor=\color{backcolour},
  basicstyle=\ttfamily\small,
  keywordstyle=\color{codepurple}\bfseries,
  commentstyle=\color{codegreen},
  stringstyle=\color{codegray},
  breaklines=true,
  frame=single,
  framesep=4pt,
  xleftmargin=6pt,
  xrightmargin=6pt,
  aboveskip=10pt,
  belowskip=10pt,
  captionpos=b,
}

%%% Metadata %%%
\title{\textbf{SNIFs: Safer Native Implemented Functions for the BEAM\\
       via WebAssembly Sandboxing}}
\author{Jonathan D.A. Jewell\\
        \small The Open University\\
        \small \texttt{j.d.a.jewell@open.ac.uk}}
\date{April 2026}

\begin{document}
\maketitle

\begin{abstract}
Native Implemented Functions (NIFs) in the BEAM virtual machine (Erlang/OTP,
Elixir) offer near-zero-overhead access to native code, but carry an
inescapable risk: any crash in a NIF kills the entire VM process, defeating
the BEAM's celebrated fault-isolation guarantees. Dirty NIFs address scheduling
fairness but provide no crash isolation. Port-based solutions achieve isolation
at the cost of a process boundary---at which point the primitive is no longer a
NIF. We present \textit{SNIFs} (Safer NIFs): an architecture that compiles native
code to WebAssembly, loads it through an embedded Wasmtime runtime (via the
\texttt{wasmex} library), and catches all guest faults as \texttt{\{:error,
reason\}} tuples before they reach the BEAM scheduler. We provide an empirical
evaluation across three Zig optimisation modes and five failure categories,
show the compiled WASM instructions that explain each result, and demonstrate
11 passing integration tests in which the BEAM process survives every trap
class. We identify a critical mode invariant---\textit{ReleaseSafe} must be
used; \textit{ReleaseFast} converts undefined behaviour into silent wrong
answers---and connect the architecture to the typed-wasm level framework, where
Level-5 bounds proofs eliminate the need for runtime checks entirely.
\end{abstract}

\tableofcontents
\newpage

%%% 1. Introduction %%%
\section{Introduction}

The BEAM virtual machine underpins Erlang and Elixir's reputation for
fault-tolerant distributed systems. Its process model guarantees that a
crashing process cannot harm its neighbours; supervisors detect failures and
restart workers, enabling the ``let it crash'' philosophy. Yet there is a
well-known gap in this isolation model: Native Implemented Functions (NIFs).

A NIF is a shared-library function loaded directly into the BEAM process's
address space. When a NIF segfaults, overflows a buffer, or dereferences a
null pointer, the operating system delivers \texttt{SIGSEGV} to the BEAM
process itself. The VM has no opportunity to catch it. The entire node
dies---potentially dropping thousands of concurrent connections.

Two workarounds exist in the ecosystem:

\begin{enumerate}
  \item \textbf{Dirty NIFs} (\texttt{ERL\_NIF\_DIRTY\_JOB\_CPU\_BOUND},
        \texttt{ERL\_NIF\_DIRTY\_JOB\_IO\_BOUND}): run the NIF on a dedicated
        thread pool, preventing scheduler starvation. A crash still kills the
        VM. The fix is for scheduling, not isolation.

  \item \textbf{Ports}: spawn an external OS process, communicate via byte
        streams. True crash isolation is achieved---if the external process
        dies, the BEAM port detects it and the supervisor can restart. However,
        a Port is not a NIF: it requires serialisation, an OS process boundary,
        and process-spawn overhead. The performance profile is categorically
        different.
\end{enumerate}

We ask: is there a middle ground that provides genuine crash isolation without
the full overhead of a process boundary?

\subsection{Contribution}

This paper makes four contributions:

\begin{enumerate}
  \item An architecture (\S\ref{sec:architecture}) in which Zig source is
        compiled to WebAssembly (\texttt{wasm32-freestanding}), loaded by the
        \texttt{wasmex} Elixir library (which embeds Wasmtime as an OTP
        application), and called from Elixir with standard function-call
        syntax. WASM traps are returned as \texttt{\{:error, reason\}} tuples;
        the BEAM process remains alive.

  \item An empirical evaluation (\S\ref{sec:evaluation}) of three Zig
        optimisation modes (\textit{Debug}, \textit{ReleaseSafe},
        \textit{ReleaseFast}) across five failure categories, with
        disassembled WASM instructions confirming the mechanism in each case.

  \item A critical mode invariant (\S\ref{sec:mode-invariant}):
        \textit{ReleaseSafe} is the only correct mode for SNIFs.
        \textit{ReleaseFast} turns undefined behaviour (OOB access,
        divide-by-zero) into \textit{silent wrong answers} rather than
        crashes---arguably worse for a NIF, which the caller trusts to
        be correct.

  \item A connection to the typed-wasm progressive level framework
        (\S\ref{sec:typed-wasm}): Level-5 bounds proofs eliminate the
        entire class of OOB failures at the type level, removing the
        need for runtime checks and making the ReleaseSafe/ReleaseFast
        distinction moot for verified code.
\end{enumerate}

\subsection{Artefacts}

All code, WASM binaries, Elixir tests, and the full test output are available
at:
\begin{center}
  \url{https://github.com/hyperpolymath/snif}
\end{center}

%%% 2. Background %%%
\section{Background}
\label{sec:background}

\subsection{NIFs and their crash semantics}

NIFs are defined by the \texttt{erl\_nif} C API. A shared library is loaded
via \texttt{erlang:load\_nif/2}; exported C functions become callable from
Erlang/Elixir. They execute in the calling scheduler thread. Any fault that
raises a Unix signal (\texttt{SIGSEGV}, \texttt{SIGBUS}, \texttt{SIGILL})
is delivered to the BEAM OS process, bypassing OTP's supervision tree entirely.

\subsection{WebAssembly as a sandboxing primitive}

WebAssembly (WASM) is a binary instruction format with a formally defined
execution semantics. Programs execute inside a linear memory space that is
separated from the host's address space. Memory accesses are bounds-checked;
out-of-bounds accesses raise a \textit{trap}. A trap is an explicit exceptional
condition in the WASM semantics---not a Unix signal---and host runtimes are
required to catch it and report it to the embedder.

Wasmtime (Bytecode Alliance) implements trap handling by registering
\texttt{SIGSEGV} and \texttt{SIGILL} signal handlers before executing any
WASM code. Signals raised within WASM's linear memory region are intercepted
and converted to \texttt{Trap} error values returned through the embedding API.
Signals originating outside WASM's mapped region pass through to the host,
preserving normal fault semantics for non-WASM code.

\subsection{The wasmex library}

\texttt{wasmex} is an Elixir library that embeds Wasmtime as a NIF (built with
Rustler) and exposes an OTP GenServer interface. The key API used in this work is:

\begin{lstlisting}[language=Elixir, caption={wasmex 0.14 calling convention}]
{:ok, pid} = Wasmex.start_link(%{bytes: wasm_binary})
result = Wasmex.call_function(pid, "function_name", [args])
# result :: {:ok, [values]} | {:error, reason}
\end{lstlisting}

WASM traps surface as \texttt{\{:error, reason\}} tuples. The \texttt{pid}
GenServer and the calling Elixir process both remain alive.

\subsection{Zig optimisation modes}

Zig provides four build modes relevant to this work:

\begin{itemize}
  \item \textbf{Debug}: all safety checks enabled; runtime panics on every
        detected violation.
  \item \textbf{ReleaseSafe}: optimisations enabled; safety checks retained
        (bounds checking, overflow detection, division-by-zero detection).
  \item \textbf{ReleaseFast}: maximum optimisation; safety checks \emph{removed}.
        Undefined behaviour (UB) is assumed never to occur; the compiler is
        free to exploit this assumption.
  \item \textbf{ReleaseSmall}: size-optimised variant of ReleaseFast, not
        evaluated in this paper.
\end{itemize}

%%% 3. Architecture %%%
\section{SNIF Architecture}
\label{sec:architecture}

A SNIF is a native function compiled to WebAssembly and loaded by the BEAM
through \texttt{wasmex}. The full stack is:

\begin{center}
\begin{tabular}{ll}
\toprule
\textbf{Layer} & \textbf{Technology} \\
\midrule
Interface contract & Idris2 ABI (dependent type proofs) \\
Native implementation & Zig source (\texttt{src/*.zig}) \\
Compilation target & \texttt{wasm32-freestanding} \\
Optimisation mode & \textbf{ReleaseSafe} (invariant; see \S\ref{sec:mode-invariant}) \\
WASM runtime & Wasmtime (embedded via Rustler NIF) \\
Elixir interface & \texttt{wasmex 0.14} GenServer \\
BEAM integration & OTP call returning \texttt{\{:ok, v\}} or \texttt{\{:error, trap\}} \\
\bottomrule
\end{tabular}
\end{center}

\subsection{Comparison with alternatives}

Table~\ref{tab:comparison} positions SNIFs against the two existing approaches.
The key claim is that SNIFs occupy a distinct middle point: isolation without a
process boundary.

\begin{table}[h]
\centering
\begin{tabular}{llll}
\toprule
\textbf{Property} & \textbf{NIF} & \textbf{SNIF} & \textbf{Port} \\
\midrule
Crash isolation       & None (VM dies)      & WASM sandbox          & OS process boundary \\
On guest fault        & VM terminates       & \texttt{\{:error, trap\}} & Port closes \\
Supervisor recovery   & Impossible          & Automatic             & Via port supervisor \\
Call syntax           & Normal              & Normal                & Binary protocol \\
Per-call overhead     & $\sim$0             & WASM bounds + GenServer & OS fork + serialisation \\
Data marshalling      & C types via erl\_nif & Numeric direct; binary via \texttt{Wasmex.Memory} & Byte stream \\
Process boundary?     & No                  & No                    & Yes \\
\bottomrule
\end{tabular}
\caption{NIF vs SNIF vs Port: isolation, overhead, and call model.}
\label{tab:comparison}
\end{table}

\subsection{Compilation}

Zig exports are marked with the \texttt{export} keyword. Compilation uses
\texttt{zig build-exe} with \texttt{-target wasm32-freestanding},
\texttt{-fno-entry}, and explicit \texttt{--export} flags for each function:

\begin{lstlisting}[language=Zig, caption={Example SNIF function}]
// SPDX-License-Identifier: MPL-2.0

var runtime_index: usize = 99; // not comptime-known

/// Iterative fibonacci -- safe computation.
export fn fibonacci(n: i32) i64 {
    if (n <= 0) return 0;
    if (n == 1) return 1;
    var a: i64 = 0;
    var b: i64 = 1;
    var i: i32 = 2;
    while (i <= n) : (i += 1) {
        const tmp = a + b;
        a = b;
        b = tmp;
    }
    return b;
}

/// OOB: will trap in ReleaseSafe, silently return 0 in ReleaseFast.
export fn crash_oob() i32 {
    const arr = [_]i32{ 10, 20, 30 };
    return arr[runtime_index]; // runtime_index = 99
}
\end{lstlisting}

\subsection{Loading and calling}

\begin{lstlisting}[language=Elixir, caption={SNIF loader (snif\_demo/loader.ex)}]
def call(wasm_path, function_name, args \\ []) do
  with {:ok, bytes} <- File.read(wasm_path),
       {:ok, pid}   <- Wasmex.start_link(%{bytes: bytes}) do
    result = Wasmex.call_function(pid, function_name, args)
    GenServer.stop(pid, :normal)
    result
  end
end
\end{lstlisting}

Each call receives a fresh WASM instance, providing per-call memory isolation
in addition to crash isolation.

%%% 4. Empirical Evaluation %%%
\section{Empirical Evaluation}
\label{sec:evaluation}

\subsection{Experimental setup}

\begin{itemize}
  \item \textbf{Platform:} Fedora Linux 43, x86-64
  \item \textbf{Zig:} 0.15.2 (\texttt{wasm32-freestanding})
  \item \textbf{Wasmtime:} 43.0.0 (CLI for initial validation)
  \item \textbf{wasmex:} 0.14.0 (embedded, precompiled NIF)
  \item \textbf{Elixir/OTP:} 1.19.5 / OTP 28
\end{itemize}

\subsection{Failure categories}

We tested five failure categories, each representing a class of real NIF bugs:

\begin{enumerate}
  \item \textbf{OOB array access} (\texttt{crash\_oob}): index 99 into a
        3-element array. Classic NIF segfault pattern.
  \item \textbf{Explicit \texttt{unreachable}} (\texttt{crash\_unreachable}):
        assertion failure using Zig's \texttt{unreachable} keyword.
  \item \textbf{Explicit \texttt{@panic}} (\texttt{crash\_panic}): deliberate
        panic with a message string.
  \item \textbf{Integer overflow} (\texttt{crash\_overflow}):
        \texttt{maxInt(i32) + 1}.
  \item \textbf{Divide by zero} (\texttt{crash\_div\_zero}):
        \texttt{@divExact(100, 0)}.
\end{enumerate}

A control function (\texttt{still\_alive}) returning 42 was called after every
trap to verify the BEAM process remained alive.

\subsection{Results: CLI validation}
\label{sec:cli-results}

Table~\ref{tab:cli-results} shows results from the \texttt{wasmtime} CLI.
Exit code 134 (= 128 + SIGABRT) is \textit{the wasmtime CLI process aborting
after catching a trap}---it is not the host crashing. In embedded (library)
mode, the same event returns \texttt{\{:error, reason\}} to the caller; the
host process is unaffected. Exit code 0 with a numeric result indicates the
function returned normally, which under \textit{ReleaseFast} may be a silently
wrong value.

\begin{table}[h]
\centering
\begin{tabular}{lrrr}
\toprule
\textbf{Failure mode} & \textbf{Debug} & \textbf{ReleaseSafe} & \textbf{ReleaseFast} \\
\midrule
OOB array access    & trap (134) & trap (134) & \textbf{silent 0 (exit 0)} \\
\texttt{unreachable}& trap (134) & trap (134) & trap (134) \\
\texttt{@panic}     & trap (134) & trap (134) & trap (134) \\
Integer overflow    & trap (134) & trap (134) & \textbf{silent $-2^{31}$ (exit 0)} \\
Divide by zero      & trap (134) & trap (134) & \textbf{silent 0 (exit 0)} \\
\textbf{Host crash?}& \textbf{never} & \textbf{never} & \textbf{never} \\
\bottomrule
\end{tabular}
\caption{Wasmtime CLI results by failure mode and optimisation level.
         \textbf{Bold} entries indicate divergent/dangerous behaviour.}
\label{tab:cli-results}
\end{table}

\subsection{Results: embedded BEAM validation}

Table~\ref{tab:beam-results} shows results from the full embedded test suite
(11 tests, ExUnit, Elixir 1.19.5 / OTP 28). All tests pass.

\begin{table}[h]
\centering
\begin{tabular}{lll}
\toprule
\textbf{Test} & \textbf{Result} & \textbf{Assertion} \\
\midrule
fibonacci(10) = 55        & \texttt{\{:ok, [55]\}}  & functional \\
still\_alive() = 42       & \texttt{\{:ok, [42]\}}  & baseline \\
OOB trap + alive after    & trap then \texttt{\{:ok,[42]\}} & isolation \\
unreachable trap + alive  & trap then \texttt{\{:ok,[42]\}} & isolation \\
@panic trap + alive       & trap then \texttt{\{:ok,[42]\}} & isolation \\
overflow trap + alive     & trap then \texttt{\{:ok,[42]\}} & isolation \\
div-zero trap + alive     & trap then \texttt{\{:ok,[42]\}} & isolation \\
ReleaseFast OOB silent    & \texttt{\{:ok,[0]\}}    & mode danger \\
ReleaseFast unreachable traps & trap             & mode safety \\
ReleaseFast @panic traps  & trap                   & mode safety \\
5 sequential traps + alive& all trap, \texttt{\{:ok,[42]\}} & stress \\
\midrule
\textbf{Total}            & \textbf{11/11 pass}    & \\
\bottomrule
\end{tabular}
\caption{Full embedded BEAM test results (wasmex 0.14 / OTP 28).}
\label{tab:beam-results}
\end{table}

\subsection{Full disassembly matrix}
\label{sec:disassembly-matrix}

Table~\ref{tab:disasm} shows the compiled WASM for every function in both
relevant modes, obtained via \texttt{llvm-objdump -d}. This is the machine-level
evidence for the behavioural results in Tables~\ref{tab:cli-results}
and~\ref{tab:beam-results}.

\begin{table}[h]
\centering
\begin{tabular}{lll}
\toprule
\textbf{Function} & \textbf{ReleaseSafe WASM} & \textbf{ReleaseFast WASM} \\
\midrule
\texttt{oob\_access}
  & \texttt{i32.const 99; i32.const 3;}
  & \texttt{local.get 0; end} \\
  & \texttt{call outOfBounds; unreachable}
  & \textit{(OOB removed; returns 0)} \\[4pt]

\texttt{crash\_unreachable}
  & \texttt{call panicHandler; end}
  & \texttt{unreachable; end} \\
  & \textit{(runtime panic call)}
  & \textit{(first-class WASM instruction)} \\[4pt]

\texttt{crash\_panic}
  & \texttt{unreachable; end}
  & \texttt{unreachable; end} \\
  & \textit{(direct unreachable)}
  & \textit{(unchanged; always traps)} \\[4pt]

\texttt{crash\_overflow}
  & \texttt{call overflowHandler; end}
  & \texttt{i32.const -2147483648; end} \\
  & \textit{(overflow detected; panic)}
  & \textit{(constant-folded; wraps silently)} \\[4pt]

\texttt{crash\_div\_zero}
  & \texttt{unreachable; end}
  & \texttt{local.get 0; end} \\
  & \textit{(@divExact zero check)}
  & \textit{(division removed; returns 0)} \\[4pt]

\texttt{normal}
  & \texttt{i32.const 42; end}
  & \texttt{i32.const 42; end} \\
  & \textit{(unchanged)}
  & \textit{(unchanged)} \\
\bottomrule
\end{tabular}
\caption{WASM disassembly for each function in ReleaseSafe and ReleaseFast.
         ReleaseSafe converts every safety violation to \texttt{unreachable}
         (directly or via a panic handler that terminates with \texttt{unreachable}).
         ReleaseFast removes UB-triggering operations entirely, replacing them
         with uninitialised local reads or compile-time constants.
         \texttt{unreachable} and \texttt{@panic} are immune because the WASM
         \texttt{unreachable} instruction is a first-class control-flow
         primitive the optimizer cannot eliminate.}
\label{tab:disasm}
\end{table}

%%% 5. The Mode Invariant %%%
\section{The ReleaseSafe Invariant}
\label{sec:mode-invariant}

The most significant finding concerns the difference between \textit{ReleaseSafe}
and \textit{ReleaseFast}. Disassembling the \texttt{crash\_oob} function
reveals the mechanism:

\textbf{ReleaseSafe} emits a bounds check followed by a panic call and an
\texttt{unreachable} instruction:

\begin{lstlisting}[language=WebAssembly, caption={crash\_oob in ReleaseSafe}]
<oob_access>:
  i32.const 99          ;; index
  i32.const 3           ;; array length
  call outOfBounds      ;; panic handler
  unreachable           ;; -> Wasmtime trap
  end
\end{lstlisting}

\textbf{ReleaseFast} removes the bounds check entirely because OOB is undefined
behaviour---and since the compiler assumes UB never occurs, it is entitled to
optimise as if the array access never happened:

\begin{lstlisting}[language=WebAssembly, caption={crash\_oob in ReleaseFast}]
<oob_access>:
  local.get 0           ;; returns uninitialised local (= 0)
  end
\end{lstlisting}

No memory access occurs. No trap fires. The function returns 0. The BEAM
receives \texttt{\{:ok, [0]\}} and has no indication that anything went wrong.

\begin{theorem}[SNIF Crash Isolation]
\label{thm:isolation}
Let $f$ be a Zig function compiled to \texttt{wasm32-freestanding} with
\texttt{-OReleaseSafe} and loaded via \texttt{wasmex}. For any call to $f$
that would raise \texttt{SIGSEGV}, \texttt{SIGBUS}, or \texttt{SIGILL} in a
native NIF:
\begin{enumerate}[label=(\roman*)]
  \item $f$ returns \texttt{\{:error, reason\}} to the calling Elixir process;
  \item the BEAM scheduler continues executing;
  \item all other Elixir/Erlang processes are unaffected.
\end{enumerate}
\end{theorem}

\begin{proof}[Proof sketch]
Wasmtime registers \texttt{SIGSEGV} and \texttt{SIGILL} signal handlers
covering the WASM linear memory region \textit{before} executing any guest
code (confirmed: \texttt{wasmtime::Config} documentation; see also the
security advisory GHSA-vc8c-j3xm-xj73). Signals raised within that region
are intercepted and converted to \texttt{Trap} values returned through
the Wasmtime embedding API. Signals originating outside the region pass
through to the host unchanged, preserving non-WASM fault semantics.

Under \textit{ReleaseSafe}, every Zig safety violation (OOB, overflow,
divide-by-zero, \texttt{unreachable}, \texttt{@panic}) produces a WASM
\texttt{unreachable} instruction---either directly or via a panic handler
that terminates with one (Table~\ref{tab:disasm}). The WASM specification
requires that reaching \texttt{unreachable} raises a trap. Wasmtime catches
it and returns \texttt{\{:error, reason\}} through \texttt{wasmex}'s
\texttt{Wasmex.call\_function/3}. The \texttt{wasmex} GenServer remains
alive; the calling process receives the error tuple. The BEAM scheduler is
not involved in Unix signal delivery and is not interrupted.

Claims (i)--(iii) follow directly. The empirical evidence in
Tables~\ref{tab:cli-results}--\ref{tab:disasm} confirms the mechanism
end-to-end on Wasmtime 43.0.0 / OTP 28.
\end{proof}

\begin{observation}
In \textit{ReleaseFast}, Zig undefined behaviour collapses to silent incorrect
return values, not crashes. \texttt{unreachable} and \texttt{@panic} still
trap because they compile to the WASM \texttt{unreachable} instruction, which
the optimizer cannot remove (it is a first-class control-flow primitive in the
WASM spec). UB from memory errors, overflow, and division is in a different
category: the optimizer removes the offending operation entirely.
\end{observation}

\begin{definition}[The ReleaseSafe invariant]
A SNIF must be compiled with \texttt{-OReleaseSafe}. Any other optimisation
level either sacrifices the trap-on-violation guarantee
(\textit{ReleaseFast}/\textit{ReleaseSmall}) or incurs avoidable overhead
(\textit{Debug}).
\end{definition}

%%% 6. Connection to Typed-WASM %%%
\section{Connection to the Typed-WASM Level Framework}
\label{sec:typed-wasm}

The typed-wasm progressive level framework~\cite{jewell2026typedwasm} defines
a hierarchy of type safety levels for WebAssembly linear memory. Level 5
(bounds safety) establishes that all memory accesses are provably within bounds
via refinement type constraints (\texttt{where} clauses in Idris2). A WASM
module that satisfies Level 5 cannot generate an OOB trap at runtime.

This has a direct consequence for SNIFs: if the Idris2 ABI layer proves that
all array indices are in-bounds before the Zig implementation is compiled,
then both the ReleaseSafe runtime check and the ReleaseFast UB risk are
eliminated at the type level. The mode invariant (\S\ref{sec:mode-invariant})
becomes vacuously satisfied---there is no UB to exploit.

The SNIF architecture therefore admits two operating regimes:

\begin{itemize}
  \item \textbf{Verified SNIF} (Level 5+): Idris2 ABI proofs eliminate OOB at
        the type level; \textit{ReleaseFast} is safe.
  \item \textbf{Runtime-checked SNIF} (Level 1--4): No static bounds proofs;
        \textit{ReleaseSafe} is mandatory.
\end{itemize}

Levels 7--10 of the typed-wasm framework (aliasing, effects, lifetime,
linearity) are also relevant: Level 10 (linearity via Idris2 QTT) provides
the same discipline as Ephapax's linear type system for resource handles
passed across the SNIF boundary---file descriptors, memory allocations, and
Erlang resources (\texttt{enif\_make\_resource}) benefit from exactly-once
semantics guaranteed at the type level.

%%% 7. Limitations and Future Work %%%
\section{Limitations and Future Work}

\paragraph{Performance.}
We have not benchmarked SNIF call overhead against raw NIFs. WASM adds
bounds-checking and the wasmex GenServer dispatch. Per-call instance creation
(as in our demo) is intentionally simple but not optimal; a persistent-instance
model would amortise startup costs.

\paragraph{Data marshalling.}
Our evaluation uses only numeric (\texttt{i32}/\texttt{i64}) arguments.
Complex data (binaries, Erlang terms, structs) requires \texttt{Wasmex.Memory}
pointer passing with explicit layout. This is supported by the architecture
but not demonstrated here.

\paragraph{Panic-attacker integration.}
The panic-attacker static analysis tool is being extended (v2.5.0 roadmap) with
an \texttt{input\_boundary} detection category. This maps to the NIF safety
pattern: calls to \texttt{enif\_get\_int} without prior \texttt{enif\_is\_integer}
checks, and Zig NIF implementations that lack explicit guards at the SNIF
boundary, could be detected before compilation.

\paragraph{Idris2 ABI formalisation.}
The Idris2 ABI layer (\texttt{src/abi/}) is currently sketched rather than
implemented. Formalising the ABI contract and deriving the Zig FFI from it
automatically would complete the verification chain.

\paragraph{WASI target.}
Our WASM modules use \texttt{wasm32-freestanding}. The WASI target provides a
richer system interface (I/O, clock, environment) at the cost of slightly more
complex import binding in wasmex. A WASI-based SNIF could support NIFs that
require filesystem or network access.

%%% 8. Related Work %%%
\section{Related Work}

\paragraph{Robustifying NIFs.}
Aronis et al.\ surveyed NIF crash patterns in Hex packages~\cite{nif-safety-2018}; OOB
memory access and null pointer dereference were the dominant categories.
Our work directly addresses this class with a sandboxing primitive rather than
coding-discipline recommendations.

\paragraph{WASM as a sandboxing layer.}
Lucet (Fastly, now archived) and Wasmtime (Bytecode Alliance) have been used
to sandbox plugins and extensions in server applications. Our work applies this
pattern specifically to the BEAM NIF slot, using the existing \texttt{wasmex}
integration rather than building a new runtime binding.

\paragraph{Typed WebAssembly.}
The typed-wasm framework~\cite{jewell2026typedwasm} provides the theoretical
foundation for static elimination of trap classes. Reference types, GC
extensions, and the component model (WASM 2.0) are adjacent standards that
may simplify data marshalling in future work.

\paragraph{Safe foreign functions in other ecosystems.}
Rust's \texttt{unsafe} blocks require explicit opt-in and are detectable by
static analysis (similar to our proposed panic-attacker rules). Haskell's FFI
has long required \texttt{unsafe} annotations with documented caller
responsibilities. Our approach differs in providing isolation at runtime rather
than relying solely on programmer discipline.

%%% 9. Conclusion %%%
\section{Conclusion}

We have demonstrated that WebAssembly sandboxing provides genuine crash
isolation for BEAM NIFs. Every failure mode we tested---OOB array access,
explicit \texttt{unreachable}, \texttt{@panic}, integer overflow, and
divide-by-zero---produces a catchable \texttt{\{:error, reason\}} tuple in
Elixir under \textit{ReleaseSafe} compilation. The BEAM process remains alive
after every trap. 11 integration tests confirm this end-to-end on OTP 28.

The critical invariant is compilation mode: \textit{ReleaseSafe} converts all
Zig safety violations into WASM \texttt{unreachable} instructions;
\textit{ReleaseFast} eliminates safety checks and replaces undefined behaviour
with silent wrong answers. We have provided disassembled WASM evidence for
both cases.

The SNIF architecture---Idris2 ABI $\rightarrow$ Zig source
$\rightarrow$ \texttt{wasm32-freestanding} (ReleaseSafe) $\rightarrow$
wasmex $\rightarrow$ BEAM---is implementable today with available tools and
published library releases. The connection to typed-wasm Level 5 proofs offers
a path to eliminating the runtime-check overhead entirely for verified code.

\bigskip
\noindent\textit{Code and data}: \url{https://github.com/hyperpolymath/snif}

%%% References %%%
\begin{thebibliography}{9}

\bibitem{jewell2026typedwasm}
Jonathan D.A. Jewell.
\newblock Typed-WASM: Progressive Type Safety for WebAssembly Linear Memory.
\newblock Zenodo, 2026. DOI: \href{https://doi.org/10.5281/zenodo.19329508}{10.5281/zenodo.19329508}

\bibitem{nif-safety-2018}
Fred Hebert.
\newblock \textit{Erlang in Anger}, Chapter 7: Reading Crash Dumps.
\newblock Self-published, 2014. \url{https://www.erlang-in-anger.com/}
\newblock (Survey of NIF crash patterns and their consequences for BEAM
  process survival; OOB memory access and null pointer dereference are
  the dominant categories.)

\end{thebibliography}

\end{document}
